\section{Introduction}

A context-dependent adapter has to decide what to store and how to query it.
A larger memory can keep more information.
Spending more computation on the stored representation is a different option: an extra query may reveal a composition that the smaller state already holds.
Neither option follows from counting parameters.
We ask whether the query can improve over several updates while the input, the context and the frozen base projection stay fixed.

The test case is a chain of associations.
A query names a starting entity, and the target is the terminal of its chain.
A reader with a suitable memory map can read a successor, form a new query from it, and continue.
High accuracy on the terminal task does not by itself show that this happens.
A learned writer can expose information about the terminal set directly, so a model can answer some questions without following the path.
Proposition 10 (Appendix~\ref{app:D}) gives a conditional construction of such a statistic, and factorizing the writer inputs removes one direct term of it without removing every global statistic.

We study a low-rank adapter whose reader updates a low-dimensional query several times over a fixed delta-rule memory (Section~\ref{sec:setting}).
The writer makes one decay-first delta update per eligible context token.
Three questions organize the paper.
First, does a larger evaluation depth improve one trained reader while its input and memory stay the same?
Second, do separately trained evolving-query systems beat single-read, fixed-query and memory-free controls at the same rank?
Third, what do these comparisons cost in parameters and measured latency?
Rank is held fixed so that extra reader computation is not mixed up with extra representation width, and extra latency is reported as a cost.  % FR1

\paragraph{What we find.}
\begin{itemize}[leftmargin=1.4em,itemsep=2pt,topsep=3pt]
\item With an oracle writer, more reads help.
In a symbolic task with orthogonal oracle codes, repeated reads improve the reachable hops and beat fixed-query and memory-free controls.  % S1 S2
In a frozen language model at rank 128, one shared module read four times beats a single read with the same parameter count on hops 2 to 3, at two backbone sizes of one family: by +21.550 pp (Qwen3-0.6B-Base) and +25.783 pp (Qwen3-1.7B-Base), in every seed.  % LP11 LP15
The four-unshared-reads reader B8 fails its registered gate at both sizes; a copy of B8 started from the shared solution keeps that solution, and the verdict stays.  % LP9 LP16 BW1
\item Learned writing is the weak point.
In the symbolic task a learned writer gives a smaller, seed-dependent benefit, and a factorized writer gives a stronger walk-specific effect but loses single-hop accuracy.  % S3 S9
In the language model, no learned writer reaches the registered 0.80 gate, including the DPFP and DPFP-X screens; each registered decision is to stop.  % LL9 DP1 DP4
A parser into a fixed table of the oracle's own codes drives the frozen reader at the oracle's level, but it is not a learned writer.  % CT1 CT2
We therefore make no learned-writer and no end-to-end claim.
\item The oracle result is bounded.
At its trained depth the reader reaches hop 3; run deeper, the same checkpoint reaches hops 4 to 8 at low accuracy while it loses hop 1.  % LP21
A label-free per-question stop raises hops 4 to 8 over the trained depth by 3.90 and 5.68 pp, but hop 2 falls by 10.83 and 12.27 pp and no efficiency claim is made.  % LS1 LS3
Per-step readouts are decoding only, and bridge patching gives no causal claim.  % LP5 PC1
\item Stopping rules, memory load and budget controls are reported as costs and limits.
Same-rank latency is a cost, and no efficiency claim is made over larger-rank single reads.  % FR1
\end{itemize}

\paragraph{Limits that apply to every claim.}
All entities come from a closed 1,024-name vocabulary, and every evaluation set draws on its seen pool.
The one exception is a descriptive scope row on 256 unseen-pool names of the same vocabulary, which supports only the statement that the depth gain is also present there.  % UN2 UN3
Nothing here concerns names outside the vocabulary or new entity codes.
Effects hold at the trained hops and depths, and longer hops appear only as measured cells.
Oracle writers and learned writers answer different questions and are never pooled, and the language-model oracle differs from the symbolic oracle.  % LP2 LP6
Every language-model composition number rests on one random oracle code table (one oracle seed) and three training seeds.  % LP22
Appendix~\ref{app:C} lists each principal claim with its comparator, writer, split, gate and evidence rows.
